% !TeX root = ../papers/nc-p3.tex
\section{等行列式の第二補完における八つの評価障害と無印商}
\label{sec:equal-determinant-second-sheet}
前節の $z=-1,\lambda=1$ の第三表示 $B=\left(\begin{smallmatrix}-d&b\\c&-a\end{smallmatrix}\right)$ を取る。
$D=ad-bc$, $p=a+d,q=a-d,r=b+c,w=b-c$ とし $D\ne0$ を保つ。
$p=0$ では $B=A$、$w=0$ では前節の第四表示に一致する。
従ってこの二平面では既に正交換と負交換の端点で全次数正則性が決定されている。
\begin{theorem}\label{thm:equal-determinant-point-configurations}
$pw\ne0$ とする。全四辺の左右の非零評価が階数3以上であるための必要十分条件は
\[
    (a^2-b^2)(d^2-c^2)(a^2-c^2)(d^2-b^2)\ne0 .
\]
条件が破れれば非標準ASである。条件が成立すれば反対階数は10、四辺はいずれも
長さ20の有限スキームの同型のグラフである。位置1の座標の平方を順に並べると、支持は
$(1,-1,0,0),(0,0,1,1)$ と次の四点の、座標ごとの平方写像による全逆像である。
\[
\begin{array}{c|cccc}
 &Y_0&Y_1&Y_2&Y_3\\ \hline
 &4b&0&(a+b)^2&(a-b)^2\\
 &0&4c&(d+c)^2&(d-c)^2\\
 &(d-c)^2&-(a-b)^2&D(d-c-a+b)&0\\
 &(d+c)^2&-(a+b)^2&0&D(a+b-c-d)
\end{array}
\]
位置3では $b,c$ を交換する（文字倍率の符号を除く）。
各位置の相異なる点の数は $20-4[b=0]-4[c=0]-4[q=w]-4[q=-w]$。
合流した各支持の局所環は $k[t]/(t^3)$、他の支持は被約である。
\end{theorem}
\begin{proof}
奇位置の座標を $y=(y_0,y_1,y_2,y_3)$、位置0の座標を $(u,v,s,t)$ とする。
厚い二色の四式を
\[
 D\binom{uy_0}{vy_1}=\begin{pmatrix}c-d&c+d\\-a+b&-a-b\end{pmatrix}\binom{sy_2}{ty_3},\quad
 D\binom{uy_1}{vy_0}=\begin{pmatrix}a+b&a-b\\c+d&c-d\end{pmatrix}\binom{sy_3}{ty_2}
\]
と取れ、薄い二式は $uy_2+vy_3+sy_0-ty_1=uy_3+vy_2+sy_1+ty_0=0$。
奇位置では全座標の独立な符号変更が点配置を保存する。実際四元群の高低二ビットの文字に加え、
前節の移動 $(1,0,1,0)$ は偶位置で両組内の交換、奇位置1で第1番目だけ、
位置3で第0番目だけの符号変更を与える（番号は0から）。
最大小行列式の各文字成分もスキーム上零となる。$y_0y_1y_2y_3$ の成分の係数に
$\pm2(b-c)D\ne0$ があるため、必ず一座標が零である。

$y_3=0,y_2=1$ なら $u=-sy_0+ty_1,v=-sy_1-ty_0$。
$L=(d-c)/D,M=(a-b)/D$ と置く。残る式は二行二列の回転
$(s,t)\mapsto(-t,s)$ の二固有線ごとに解ける。$t=\alpha s$, $\alpha^2=-1$ の線では
\[
 y_0(y_0-\alpha y_1)=L,\qquad y_1(y_0-\alpha y_1)=-\alpha M .
\]
これは四つの厚い式を $(s,t)$ と $(-t,s)$ の係数で読むことでも従う。
$L,M$ はともには零でない。$L=M$ なら解はなく、異なれば第三行の平方値が得られる。
$L,M$ の一方だけが零なら同じ $y$ で両固有線が核に入り、評価階数が2以下になる。
$y_2=0,y_3=1$ でも $L,M$ を $-(d+c)/D,-(a+b)/D$ とすれば同じ消去で第四行を得る。
この二平面の検査で $a=\pm b$ と $d=\pm c$ を全て排除した。

$y_1=0,y_0=1$, $y_2y_3\ne0$ で、上で除いた平面を避けるとする。
二番目と三番目の厚い式から $t=\alpha s$, $\alpha^2=1$、
$y_2=-\alpha(a+b)y_3/(a-b)$。$s=t=0$ は他の二式により不可である。
残る二つの厚い式より $u/s=2\alpha y_3/(a-b), v=\alpha u$。
薄い式から $4b y_3^2=(a-b)^2$ となる。これで第一行の四点を得る。
$b=0$ ならこのチャートに点はない。$y_0=0$ 側は同じ代入で第二行となる。
片半分だけが非零なら、厚い式により対応する半分の相手座標が零となり、
薄い二式から低位の平方比 $-1$ または高位の平方比 $1$ の二点となる。
いずれも核は直線である。各三座標の点でも記載の比と非零の分母から核は直線である。
二位置移動は二つの角を転置するので、もう一方の奇位置で同じ判定をすると
$a=\pm c,d=\pm b$ が加わる。片側反転の評価下界定理\ref{thm:single-opposite-evaluation-ranks}
より、八平面のいずれも非標準ASである。

残った全点で奇位置の射影は有限で核が直線だから、辺自体も有限である。
従って偶位置の核も高々直線である。点対は六つの双一次式の期待余次元の完全交叉で長さ20であり、
両射影は最大小行列式のスキームへの同型である。
反対切断では色0の行列式が $8p^2/D$、薄い二色は各階数2。
色1の二組の行は組ごとに等しく、残る二行は $(b,c,1,1)$,
$(1,-1,(-q+w)/D,(q+w)/D)$ である。比例すれば $q=r=0$ かつ $a^2=b^2$ で除かれている。
従って階数の和は10である。

点対の族は上の評価開集合で有限平坦である。実際、滑らかな基底上の相対完全交叉で
全ファイバーが同じ零次元を持つ。一般配置からの特殊化で第一・第二行の各四点は
$b=0,c=0$ でそれぞれ純高位二点に二点ずつ合流する。同時の消滅は $w\ne0$ に反する。
第三・第四行も $q=w,q=-w$ で純低位二点に二点ずつ合流し、この二条件も同時には起きない。
他に重なりはない。局所長さ3のGorenstein環は接次元1であるから局所構造も従う。
元の六式、符号成分と後の置換作用は \path{computations/verify_equal_determinant_second_sheet.py} が検算する。
\end{proof}

\begin{theorem}\label{thm:equal-determinant-generic-quotient}
$x=q/p,y=r/p,s=w/p$ とする。射影安定群の位数8の部分では、指定始点の無印同型は
\[
    Q=x^2+y^2,\qquad L=xys,\qquad H=s^2
\]
の一致と同値である。この軌跡は稠密開集合を含み、その商は有理三次元である。
一位置だけの周期は反対階数10の正則点では起きず、二位置移動は $L$ の符号を変える。
従ってこの部分の正則点では $L=0$ で周期2、他で周期4である。
評価開集合の残りにも、支持点の有限枠と文字倍率による無印同型・周期の必要十分判定がある。
\end{theorem}
\begin{proof}
四元群の二指標を $H_0$（高位ビット）と $J$（低位ビット）、前記の追加対称を $S_0$ とする。
射影的には可換な三つの対合である。偶位置の線型持上げでは $J$ と $S_0$ だけが反可換となる。
従ってこの群のうち四位置の文字線を与え得る階数2の群は、交換子形式の根基
$\langle H_0\rangle$ を含まねばならない。$\langle H_0,S_0\rangle$ は奇位置1の
高位二文字が同じ文字を持ち、$\langle H_0,JS_0\rangle$ では低位二文字が同じになる。
よって四位置で正則表現を与える四元群は元のもの一つである。
系\ref{cor:minus-one-second-completion-periods}の二十点標本により一般点の安定群はちょうど位数8である。
有限エタールな点配置が枠を持つ開部分では自己同型候補の数は局所的に有限であり、
標本で失敗する候補を除けば共通の稠密開集合を取れる。

この固有な四元群を運ぶ写像は、各位置の文字基底の置換と対角倍率である。
薄い48語から線型部分は共通で、移動量の上位ビットも共通になる。
32候補から薄い符号を合わせられる指定始点の16個を取り、四角の共通倍率を除いて
$(x,y,s)$ を読むと、作用は $x,y$ の交換と
$(x,y,s)\mapsto(\epsilon x,\eta y,\epsilon\eta s)$ で尽きる。
実際符号補正のSmith方程式を解くと、$[a:b:c:d]$ に
$(a,-b,-c,d),(d,c,b,a),(p+r,q+w,q-w,p-r)$ の作用がある。
各射影変換を二置換ずつが与える。$H\ne0$ だから $x^2y^2=L^2/H$ であり記載の三数が軌道を分ける。
$D\ne0$ は $1-Q+H\ne0$。二位置移動は角の転置により $s$ の符号を変える。
周期1の排除は前節の二交換の議論による。

安定群の拡大点でも、両奇位置では表の支持が、両偶位置ではその評価核の支持が得られる。
各々が独立四点を含む。実際、二組の純半分の点の相手は反対の半分にある相異なる二点である。
全位置で独立四点の順序付き組を列挙し、その各枠で四窓の非零台と飽和整数文字関係の単項式比を取れば
共通の正規形を持つかが判定できる。局所環を点数だけで同一視する必要はなく、比較するのは元の四窓自体である。
各移動について同じ有限検査をすれば周期も決まる。これらの同型判定は正則点では添字保持同型を与える。
\end{proof}
